%Chargement des paquets
\usepackage{amsmath}
\usepackage{amsfonts}
\usepackage{amssymb}
\usepackage{enumerate}
\usepackage{mathtools}
\usepackage{bbm}
\usepackage{xparse, etoolbox}
\usepackage{enumerate}
\usepackage{enumitem}
\usepackage{mathabx}
\usepackage{babel}[french]

%environment
\usepackage{amsthm}
\usepackage{keytheorems}
\usepackage{hyperref} 

%Interface théorème
\newkeytheorem{lem}[
	name = Lemme
]

\newkeytheorem{prop}[
	name = Proposition
]

\newkeytheorem{thm}[
	name = Théorème
]

\newkeytheorem{coro}[
	name = Corollaire
]

\renewcommand*{\proofname}{Démonstration}

\theoremstyle{definition}
\newtheorem{defn}{Définition}

\theoremstyle{definition}
\newtheorem*{nota}{Notation}

\theoremstyle{definition}
\newtheorem*{limi}{Liminaire}

\theoremstyle{remark}
\newtheorem{rmq}{Remarque}

\theoremstyle{plain}
\newtheorem*{fait}{Fait}


%Ensembles de nombres
\newcommand{\N} {\mathbb{N}}
\newcommand{\Ne}{\N^\ast}
\newcommand{\Z} {\mathbb{Z}}
\newcommand{\Q} {\mathbb{Q}}
\newcommand{\R} {\mathbb{R}}
\newcommand{\Rb}{\overline{\mathbb{R}}}
\newcommand{\Rp}{\R_+}
\newcommand{\Rm}{\R_-}
\newcommand{\K} {\mathbb{K}}
\newcommand{\Ket} {\mathbb{K^{\ast}}}
\newcommand{\Cx}{\mathbb{C}}

%Ensembles, tribus
\newcommand{\A}{\mathbb{A}}
\renewcommand{\P}{\mathcal{P}}
\newcommand{\T}{\mathcal{T}}
\newcommand{\F}{\mathcal{F}}
\newcommand{\C} {\mathcal{C}}
\newcommand{\ouv}{\mathcal{O}}
\newcommand{\B}{\mathcal{B}}
\newcommand{\M}{\mathcal{M}}
\newcommand{\etag}{\mathcal{E}}
\newcommand{\etagOp}{\etag^{+}(\Omega)}
\newcommand{\etagO}{\etag(\Omega)}
%%Lp avant quotientage
\newcommand{\Lic}{\mathcal{L}^1}
\newcommand{\Lpc}{\mathcal{L}^p}
\newcommand{\Linfc}{\mathcal{L}^{\infty}}
\newcommand{\Lifull}{\Lic(\Omega, \B, m)}
\newcommand{\Lpfull}{\Lpc(\Omega, \B, m)}
\newcommand{\Linffull}{\Linfc(\Omega, \B, m)}
%%Lp après quotientage
\newcommand{\Li}{L^1}
\newcommand{\Ld}{L^2}
\newcommand{\Lp}{L^p}
\newcommand{\Linf}{L^{\infty}}
\newcommand{\Listd}{\Li(\Omega, m)} 
\newcommand{\Ldstd}{\Ld(\Omega, m)} 
\newcommand{\Lpstd}{\Lp(\Omega, m)} 

%Quelques opérateurs
\DeclareMathOperator{\codim}{codim}
\DeclareMathOperator{\rg}{rg}
\DeclareMathOperator{\tr}{tr}
\DeclareMathOperator{\vect}{Vect}
\DeclareMathOperator{\ima}{Im}
\DeclareMathOperator{\id}{Id}
\DeclareMathOperator{\sign}{sign}
\DeclareMathOperator{\ext}{ext}
\newcommand{\ind}{\mathbbm{1}}
\newcommand*{\spr}[2]{\left(\,#1\,\middle|\,#2\,\right)}
\newcommand{\interior}[1]{%
  {\kern0pt#1}^{\mathrm{o}}%
}
\DeclareMathOperator{\card}{card}
\DeclareMathOperator{\ppcm}{ppcm}

%Commandes de pure flemme
\newcommand{\veps}{\varepsilon}
\newcommand{\intab}{\int_{a}^{b}}
\newcommand{\intR}{\int_{-\infty}^{\infty}}
\newcommand{\intinf}{\int_{- \infty}^{+ \infty}}
\newcommand{\equi}{\Leftrightarrow}
\newcommand{\Kev}{$\K$-espace vectoriel }
\newcommand{\Kevs}{$\K$-espaces vectoriels }
\newcommand{\Rev}{$\R$-espace vectoriel }
\newcommand{\Revs}{$\R$-espaces vectoriels }
\newcommand{\Cev}{$\Cx$-espace vectoriel }
\newcommand{\Cevs}{$\Cx$-espaces vectoriels }
\newcommand{\evn}{(E, \norm{.})}
\newcommand{\interf}[2]{[#1,#2]}
\newcommand{\limb}{\lim\limits_}
\newcommand{\lebn}{\lambda_n}
\newcommand{\pp}{\text{~presque partout}}
\newcommand{\derivp}[2]{\frac{\partial #1}{\partial #2}}
\newcommand{\famiu}{(u_i)_{i \in I}}
\newcommand{\sev}{\text{~s.e.v.}}
\newcommand{\Mnp}{M_{n,p}(\K)}
\newcommand{\Mn}{M_{n}(\K)}
\newcommand{\tq}{\text{~t.q.~}}
\newcommand{\mq}{~montrer~que~}
\newcommand{\et}{\quad \text{et} \quad}
\newcommand{\ou}{\text{~ou~}}
\newcommand{\Kx}{\K[X]}


%Commandes de gars chelou
\renewcommand{\subset}{\subseteq}

%Commande restriction, ty duckduckgo
\def\restr#1#2{\mathchoice
              {\setbox1\hbox{${\displaystyle #1}_{\scriptstyle #2}$}
              \restraux{#1}{#2}}
              {\setbox1\hbox{${\textstyle #1}_{\scriptstyle #2}$}
              \restraux{#1}{#2}}
              {\setbox1\hbox{${\scriptstyle #1}_{\scriptscriptstyle #2}$}
              \restraux{#1}{#2}}
              {\setbox1\hbox{${\scriptscriptstyle #1}_{\scriptscriptstyle #2}$}
              \restraux{#1}{#2}}}
\def\restraux#1#2{{#1\,\smash{\vrule height .8\ht1 depth .85\dp1}}_{\,#2}} 

%Quelques normes
\DeclarePairedDelimiter\abs{\lvert}{\rvert}
\DeclarePairedDelimiter\labs{\bigl\lvert}{\bigr\rvert}
\DeclarePairedDelimiter\norm{\lVert}{\rVert}
\DeclarePairedDelimiterXPP\normi[1]{}\lVert\rVert{_1}{\ifblank{#1}{\:\cdot\:}{#1}}
\DeclarePairedDelimiterXPP\normd[1]{}\lVert\rVert{_2}{\ifblank{#1}{\:\cdot\:}{#1}}
\DeclarePairedDelimiterXPP\normp[1]{}\lVert\rVert{_p}{\ifblank{#1}{\:\cdot\:}{#1}}
\DeclarePairedDelimiterXPP\normq[1]{}\lVert\rVert{_q}{\ifblank{#1}{\:\cdot\:}{#1}}
\DeclarePairedDelimiterXPP\norminf[1]{}\lVert\rVert{_\infty}{\ifblank{#1}{\:\cdot\:}{#1}}

%Bracket en angle
\DeclarePairedDelimiter\angl{\langle}{\rangle}

%Set, parenthèses
\newcommand{\set}[1]{\{ #1 \}}
\newcommand{\bigset}[1]{\bigl\{ #1 \bigr\}}
\newcommand{\Bigset}[1]{\Bigl\{ #1 \Bigr\}}
\newcommand{\bigpar}[1]{\bigl( #1 \bigr)}
\newcommand{\Bigpar}[1]{\Bigl( #1 \Bigr)}

%Opitmisation des opérateurs de comparaison
\DeclarePairedDelimiter\tio{o (}{)}
\DeclarePairedDelimiter\gro{O (}{)}


\pagestyle{fancy}

\fancyhead[C]{Les théorèmes de Dini}
\fancyhead[R]{}

\begin{document}

\underline{Source :} Rombaldi - Analayse - page 239 (premier théorème) puis page 304. \newline

\underline{Leçon 241 :} Suites et séries de fonctions. Exemples et contre-exemples. \newline


\begin{thm}[Premier théorème de Dini]
\label{thm:Dini}
Soient $K$ un compact de $E, (f_n)_{n \in \N}$ une suite croissante dans $\C^0(K, \R)$ 
(i.e. pour tout $x \in K$, la suite réelle $(f_n(x))_{n \in \N}$ est croissante) et $f \in \C^0(K, \R)$. 
La suite $(f_n)_{n \in \N}$ converge uniformément vers la fonction $f$ si, et seulement si, elle converge simplement vers $f$.
\end{thm}

\begin{proof}
La condition nécessaire est évidente. Il suffit donc de démontrer la condition suffisante. \newline

%Comme $K$ est compact, il vérifie la propriété de Borel-Lebesgue. Soit $\veps >0$, pour $x \in K$, notons $\alpha_x >0$ le réel tel que :
%\[ y \in B(x, \alpha_x) \Rightarrow d(f(x), f(y)) < \veps/3.  \]
%$\alpha_x$ existe par continuité de la fonction $f$ sur K. On extrait du recouvrement :
%\[ K \subset \bigcup_{x \in K} B(x, \alpha_x) \]
%un recouvrement fini. Notons $N \geq 1$ le cardinal du recouvrement, $x_1, \dots, x_N$ les centres des boules et $\alpha_1, \dots,  \alpha_N$ les rayons associés. Pour $x \in K$ quelconque, il existe $i \in \{1, \dots, N\} \tq x \in B(x_i, \alpha_i)$ et, par suite :
%\[ d(f(x), f(x_i)) < \veps/3 . \]
%Par convergence simple de la suite $(f_n)_{n \in \N}$, on a 
%\[ \forall i \in \{1, \dots, N\}, \exists M_i \in \N, \forall n \in N : n \geq M_i \Rightarrow d(f_n(x_i),f(x_i)) < \veps/3 . \]
%Notons $M = \max \{M_i, i \in \{1, \dots, N\} \}$, pour $n \geq M$ on a 
%\[ \forall i \in \{1, \dots, N\},  d(f_n(x_i),f(x_i)) < \veps/3 . \] 
%De plus, par croissance et convergence de la suite $(f_n)_{n \in \N}$, il existe $M' \in \N$ tel que 
%\[ n \geq M' d(f(x), f_n(x)) < \veps/3 . \] 
%Finalement, en notant $M''=max(M, M')$ 
%\[ n \geq M'' \Rightarrow \forall x \in K, d(f(x), f_n(x)) < \veps, \]
%ce qui démontre bien la convergence uniforme.

Soit $\veps >0$, onsidérons les ensembles, 
\[ O_{n, \veps} = \{x \in K ; f(x)-f_n(x) < \veps \} . \] 
On remarque que, pour tout $n \in \N, O_{n, \veps}=(f-f_n)^{-1}(]-\infty, \veps[)$, 
et comme les fonctions $f-f_n$ sont continues sur $K$ pour tout $n \in \N$, 
on en déduit que les ensembles $O_{n, \veps}$ sont des ouverts.
De plus, la suite $(f_n)_\N$ étant croissante, on a $f(x) - f_{n+1}(x) \leq f(x)-f_n(x)$ donc $O_{n, \veps} \subset O_{n+1, \veps}$.
Enfin, par convergence simple de la suite $(f_n)_{n \in \N}$ vers f, on a :
\[ K \subset \bigcup_{n \in \N} O_{n, \veps} , \]
en effet, pour tout $x \in K, f(x) = \lim_\N f_n(x)$, donc il existe $n \in \N$ tel que $f(x)-f_n(x) < \veps$. \\
Par compacité de $K$, on extrait un sous-recouvrement fini de ce recouvrement, il existe donc $m \in \N$ tel que :
\[ K \subset \bigcup_{i=1}^m O_{n_i, \veps} \text{ en choisissant la suite des indices $(n_i)$ croissante,} \]
soit, par croissance des ensembles $O_{n_i, \veps}$ :
\[ K \subset O_{n_m, \veps} . \]
On conclut donc que :
\[ \exists n_m \in \N, \forall n \in \N : n \geq n_m \Rightarrow \forall x \in K, \abs{f(x)-f_n(x)} = f(x) - f_n(x) < \veps , \]
ce qui démontre la convergence uniforme de la suite $(f_n)$.
\end{proof}

\begin{thm}[Deuxième théorème de Dini]
\label{thm:Dini2}
Soit $(f_n)_\N$ une suite de fonctions continues et monotones de $[a,b]$ dans $\R$ et qui converge simplement sur $[a,b]$
vers une fonction continue $f$. Alors la convergence est uniforme.
\end{thm}

\begin{proof}
Quitte à remplacer la suite $(f_n)_\N$ par la suite $(-f_n)_\N$ on peut supposer que chaque fonction $f_n$ est croissante.
La fonction $f$ étant continue sur le compact $[a,b]$, elle est, d'après le théorème de Heine, uniformément continue sur ce compact.
Pour tout réel $\veps >0$, on peut donc exhiber $\eta >0$ tel que :
\[ \forall x,y \in [a,b], \quad \abs{x-y} < \nu \Rightarrow \abs{f(x)-f(y)} < \veps . \]
Soit $a=a_0 < \dots < a_r=b, r \geq 0$, une subdivision de $[a,b]$ de pas maximal $\eta$. 
Tout réel $x \in [a,b]$ est dans un intervalle de la forme $[a_k, a_{k+1}]$, on a donc :
\begin{align*}
\abs{f(x)-f_n(x)} & \leq \abs{f(x) - f(a_k)} + \abs{f(a_k)-f_n(a_k)} + \abs{f_n(a_k) - f_n(x)} \\
& \leq \veps + \abs{f(a_k)-f_n(a_k)} + f_n(x) - f_n(a_k) \\
& \leq \veps + \abs{f(a_k)-f_n(a_k)} + f_n(a_{k+1}) - f_n(a_k) .
\end{align*}
Or, $\lim_{n \to \infty} f_n(a_k) = f(a_k)$, il existe donc un rang $N_1 \in \N$ à partir duquel $\abs{f(a_k)-f_n(a_k)} < \veps$.
De plus, $\lim_{n \to \infty} f_n(a_{k+1}) - f_n(a_k) = f(a_{k+1}) - f(a_k) < \veps$, on peut donc trouver un second rang $N_2 \in \N$
à partir duquel $f_n(a_{k+1}) - f_n(a_k) < 2 \veps$. 
Finalement, pour $n \geq \max(N_1, N_2)$ on a :
\[ \abs{f(x)-f_n(x)} < 4 \veps , \]
et ce pour tout $x \in [a,b]$, ce qui démontre la convergence uniforme de la suite $(f_n)_\N$.
\end{proof}

\begin{appl}
La suite de fonctions $\left ( \left (1+\dfrac{x}{n} \right)^n \right)_{\Ne}$ converge uniformément vers la fonction $x \mapsto e^x$ sur
tout compact de la forme $[-a, a], a \in \R$.
\end{appl}

\begin{proof}
Les fonctions $x \mapsto \left (1+\dfrac{x}{n} \right)^n$ sont continues et croissantes sur $[-a,a]$.
De plus, comme :
\[ \left (1+\dfrac{x}{n} \right)^n = \exp \left ( n \ln \left ( 1 + \dfrac{x}{n} \right) \right ) \sim_{\infty} \exp(x), \]
on a convergence simple de la suite vers la fonction exponentielle. On déduit du théorème de de Dini $\ref{thm:Dini2}$ que la
convergence est uniforme.
\end{proof}

\end{document}